\documentclass[aspectratio=169,10pt]{beamer}\usetheme{metropolis}\title{Q08: Scaling Limits and Forecasts}\subtitle{What a fitted curve can actually tell us}\author{DLS 2026}\date{October 2026}
\begin{document}\maketitle
\begin{frame}{The question}\begin{block}{Q08}Will performance saturate or models forget as parameters grow? Can we predict more-parameter effects before training?\end{block}\end{frame}
\begin{frame}{Scaling laws are conditional models}\[L(N)\approx L_\infty+aN^{-\alpha}\]\begin{itemize}\item Fit observed loss/resource relations.\item Useful within the measured recipe and regime.\item Not a universal law of capability.\end{itemize}\end{frame}
\begin{frame}{``More parameters'' is incomplete}\begin{itemize}\item Parameters $N$, data $D$, and compute $C$ interact.\item A bigger model trained on too few tokens may be inefficient.\item Architecture, data quality, and objective also matter.\end{itemize}\vfill\small Kaplan et al. (2020); Hoffmann et al. (2022)\end{frame}
\begin{frame}{Forecast protocol}\begin{enumerate}\item Run matched small-scale experiments.\item Hold out a scale for validation.\item Fit uncertainty, not only a line.\item Declare metric and extrapolation range.\item Check a separate downstream metric.\end{enumerate}\end{frame}
\begin{frame}{Saturation has several meanings}\begin{itemize}\item Diminishing validation-loss improvement\item Benchmark ceiling or contamination\item Data-quality limit\item Inference/deployment budget limit\end{itemize}\alert{Diagnose the limiting mechanism before choosing an intervention.}\end{frame}
\begin{frame}{Forgetting also has several meanings}\begin{itemize}\item Catastrophic forgetting after task shift\item Example-level forgetting during training\item Memorization of rare or duplicated data\end{itemize}\alert{None follows automatically from parameter count.}\end{frame}
\begin{frame}{Toy experiment}\begin{itemize}\item Fit a power law to a stable synthetic regime.\item Check an interpolation point.\item Introduce a regime shift and inspect extrapolation error.\end{itemize}\vfill A fit is evidence about its data-generating regime.\end{frame}
\begin{frame}{Takeaway}\begin{block}{Use scaling laws as budgets, not oracles}State the allocation, measured metric, range, uncertainty, and regime assumptions. Evaluate downstream behavior separately.\end{block}\end{frame}\end{document}
